\documentclass[10pt]{article}

\usepackage{mathrsfs}

%\usepackage{amssymb,amsthm,amsmath,amssymb,latexsym,amsfonts,txfonts, pxfonts,wasysym,stmaryrd}
\textheight=230mm \textwidth=140mm
\topmargin=-1.5cm
\oddsidemargin=1.2cm
\evensidemargin=1.2cm
\setlength\parindent{0pt}
\parskip=7pt
%\setcounter{page}{1}
%\setlength{\baselineskip}{24pt} \baselineskip=24pt
\usepackage{amsmath,amsthm,amssymb}
\usepackage{caption}
\usepackage{graphicx}
\usepackage{bbm}
\usepackage{xcolor}
\usepackage[utf8]{inputenc}
\usepackage[english]{babel}
\usepackage{enumitem}
\usepackage{hyperref}
\usepackage{color}
\usepackage{colortbl}
\usepackage{tikz}
\usepackage{subfig}
\usepackage{wrapfig}
\usepackage[authoryear,square]{natbib}
%\usepackage{vector}

\definecolor{gblue}{RGB}{66,133,244}
\definecolor{gred}{RGB}{234, 67, 53}
\definecolor{ggreen}{RGB}{52, 168, 83}
\definecolor{gyellow}{RGB}{251, 188, 5}
\definecolor{nupurple}{RGB}{078,042, 132}
\hypersetup{
    colorlinks=true,
    linkcolor=gblue,
    filecolor=gyellow,      
    urlcolor=gred,
    citecolor=nupurple,
    % bookmarks=true,
}


\font\tencmmib=cmmib10 \skewchar\tencmmib '60
\newfam\cmmibfam
\textfont\cmmibfam=\tencmmib
\def\Ph{\mathchar'4010}
\def\Ps{\mathchar'4011}
\font\tensmc=cmcsc10
\def\smc{\tensmc}
\def\srm{\sevenrm }
\font\tenmsb=msbm10
\font\eightmsb=msbm10 scaled 1100
\def\Bbb#1{\hbox{\tenmsb#1}}
\def\Bbbb#1{\hbox{\eightmsb#1}}
\def\bbox{\quad\hbox{\vrule \vbox{\hrule \vskip2pt \hbox{\hskip2pt
\vbox{\hsize=1pt}\hskip2pt} \vskip2pt\hrule}\vrule}}
\def\lessim{\ \lower4pt\hbox{$
\buildrel{\displaystyle <}\over\sim$}\ }
\def\gessim{\ \lower4pt\hbox{$\buildrel{\displaystyle >}
\over\sim$}\ }
\def\n{\noindent}
\def\P{{\cal P}}
\def\si{\bar{\sigma}}
\def\E{{\cal E}}
\def\LL{{\cal L}}
\def\FF{{\cal F}}
\def\SS{{\cal S}}
\def\C{{\cal C}}
\def\D{{\cal D}}
\def\H{{\cal H}}
\def\X{{\cal X}}
\def\Y{{\cal Y}}
\def\A{{\cal A}}
\def\B{{\cal B}}
\def\O{{\cal O}}
\def\M{{\cal M}}
\def\I{{\cal I}}
\def\eps{{\varepsilon}}
\def\ch{{\mbox{\rm ch}\hspace{0.4mm}}}
\def\sh{{\mbox{sh}}}
\def\th{{\mbox{th}}}
\def\Av{{\mbox{\rm{Av}}}}
\def\n{\noindent}
\def\bbe{\vskip .15pc\hangindent=.5pc\hangafter=1}
\def\goin{\to\infty}
\def\qed{\hfill\break\rightline{$\bbox$}}
\DeclareMathOperator*{\argmax}{arg\,max}
\DeclareMathOperator*{\argmin}{arg\,min}


\newcommand{\e}{\mathbb{E}}
\newcommand{\p}{\mathbb{P}}
\newcommand{\lra}{\leftrightarrow}
\newcommand{\nlra}{\nleftrightarrow}
\newcommand{\bet}{\begin{theorem}}
\newcommand{\ent}{\end{theorem}}
\newcommand{\Var}{\mathrm{Var}}
\newcommand{\g}{\Gamma}
\newcommand{\de}{\delta}
\newcommand{\norm}[1]{\left\lVert#1\right\rVert}
\newcommand{\gy}[1]{\textbf{\color{gred}[Yi: #1]}}
\newcommand{\olsi}[1]{\,\overline{\!{#1}}}
\newcommand\y{\cellcolor{blue!20}}
\newcommand\x{\cellcolor{red!20}}
\DeclareMathOperator{\Tr}{Tr}

\newtheorem{lemma}{\bf Lemma}
\newtheorem{claim}{\bf Claim}
\newtheorem{definition}{\bf Definition}
\newtheorem{theorem}{\bf Theorem}
\newtheorem*{theorem*}{\bf Theorem}
\newtheorem*{remark*}{\bf Remark}
\newtheorem{corollary}{\bf Corollary}
\newtheorem{remark}{\bf Remark}
\newtheorem{example}{\bf Example}
\newtheorem{exercise}{\bf Exercise}
\newtheorem{proposition}{\bf Proposition}
\newtheorem{question}{\bf Question}
\newtheorem{condition}{\bf Condition}
\newtheorem{acknowledgements}{\bf Acknowledgements}

\newenvironment{Proof of lemma}{\noindent{\bf Proof of Lemma}}{\hfill$\Box$\newline}
\newenvironment{Proof of claim}{\noindent{\bf Proof of claim}}{\hfill$\Box$\newline}
\newenvironment{Proof of theorem}{\noindent{\bf Proof of Theorem}}{\hfill\scriptsize{$\Box$}\newline}
\newenvironment{Proof of theorems}{\noindent{\bf Proof of Theorems}}{\hfill$\Box$\newline}
\newenvironment{Proof of proposition}{\noindent{\bf Proof of Proposition}}{\hfill$\Box$\newline}
\newenvironment{Proof of propositions}{\noindent{\bf Proof of Propositions}}{\hfill$\Box$\newline}
\newenvironment{Proof of exercise}{\noindent{\it Proof of Exercise:}}{\hfill$\Box$}
\newenvironment{Acknowledgements}{\noindent{\bf Acknowledgements}}

% To print out the paper version, please comment out the next three lines.
\definecolor{dimmedwhite}{rgb}{0.85,0.85,0.85}
% \pagecolor{dimmedwhite}
% \color{dimmedwhite}


\numberwithin{equation}{section}
\author{Yi Gu}
\newcommand{\ki}{\mathbf{Ker}_{\mathbf{I}}}
\newcommand{\ko}{\mathbf{Ker}_{\mathbf{O}}}

\setcounter{tocdepth}{1}

\begin{document}

\title{Equilibrium Distribution for t-Distributed Stochastic Neighbor Embedding with Generalized Kernels}
\maketitle

\begin{abstract}
    T-distributed stochastic neighbor embedding (t-SNE) is a well-known algorithm for visualizing high-dimensional data by finding low-dimensional representations (\cite{key}). In this paper, we study the convergence of t-SNE with generalized kernels and extend the results of \cite{auffinger-tsne}. Our work starts by giving a concrete formulation of generalized input and output kernels. Then we prove that under certain conditions, the t-SNE algorithm converges to an equilibrium distribution for a wide range of input and output kernels as the number of data points diverges.
\end{abstract}

\section{The Concept of Input and Output Kernels} \label{sec_intro}
T-distributed stochastic neighbor embedding (t-SNE) is a powerful technique for
visualizing high-dimensional data in a low-dimensional space, developed in 2008
by Laurens van der Maaten and Geoffrey Hinton (\cite{key}). Despite its wide
application (in biology for example \cite{LindermanGeorgeC.2019Fitf}, \cite{10.1093/bioinformatics/btp035}), the
theoretical aspects and mathematical rigor of t-SNE have rarely been discussed
thoroughly. Recently, Auffinger and Fletcher proved that if data is sampled
independently from a compact probability measure $\mu_X$, then the output of
t-SNE converges to an equilibrium measure that is also compact as the number of
data points diverges(\cite{auffinger-tsne}). This, together with
\cite{arora2018analysistsnealgorithmdata} and
\cite{linderman2017clusteringtsneprovably}, are some of the very first results
that rigorously establishes the theoretical guarantees of t-SNE.

We aim to generalize the results from \cite{auffinger-tsne}. In the original
version of the t-SNE algorithm, for each pair of input or output data points, a
weight is assigned based on the distance between them. This weight is used to
construct probability distributions that will be used in the loss function. A
natural question to ask is: can we generalize the result for different weights
such that the convergence still holds and the equilibrium measure still exists?
A generalized formulation allows us to further reveal the theoretical nature of
this algorithm, greatly enhancing our understanding of t-SNE and dimension
reduction. Our work starts by generalizing the algorithm with the concept of
input and output kernels.

Let $d,s \geq 2$ with $d>s$, and $n \geq 1$. Let $X_1, \cdots, X_n \in
    \mathbb{R}^d$ be the input data points and $Y_1, \cdots Y_n \in \mathbb{R}^s$
the output data points. We define the input and output kernels as follows:
\begin{definition} [Generalized Input Kernel] \label{def:input_kernel}
    The input kernel $\ki(x,y,\sigma): \mathbb{R}^d  \times \mathbb{R}^d \to \mathbb{R}^+$ is a function of the input data points and parameter $\sigma$ such that for any $x,y \in \mathbb{R}^d$ and $\sigma > 0$,
    \begin{align}
        \ki(x,y,\sigma)=\ki(y,x,\sigma).
    \end{align}
\end{definition}
\begin{definition} [Generalized Output Kernel] \label{def:output_kernel}
    The output kernel $\ko(z,w):\mathbb{R}^s \times \mathbb{R}^s \to \mathbb{R}^+$ is a function of the output data points such that for any $z,w \in \mathbb{R}^s$,
    \begin{align}
        \ko(z,w)=\ko(w,z).
    \end{align}
\end{definition}
Clearly, by letting
\begin{align}
    \ki(x,y,\sigma)=\exp\left(-\frac{||x-y||^2}{2\sigma^2}\right), \quad \ko(z,w)=\frac{1}{1+||z-w||^2},
\end{align}
we recover the original t-SNE algorithm by using both $\ki$ and $\ko$ as the weights for the input and output data points. We can also say that the original t-SNE algorithm has a Gaussian input kernel and a t-distributed output kernel.

From now on, we will use the terminology \underline{\textbf{\textit{generalized
            t-SNE}}} when we use the generalized formulation for input and output kernels
to distinguish it from the original algorithm.

Clearly, not all choices of $\ki$ and $\ko$ will yield a valid generalized
t-SNE algorithm, such as $\ki$ that diverges as $||x-y|| \to \infty$. To prove
similar results as in \cite{auffinger-tsne}, we need to impose some conditions
on both $\ki$ and $\ko$. The conditions should provide just enough constraints
for the convergence to hold, while applying to a wide range of function
choices. Towards this end, we use the following conditions for the rest of our
work. We will provide a detailed discussion on the generality of these
conditions in section \ref{tsne:sec:generality_conditions}.

\begin{condition}[Input Kernel Conditions] \label{input_kernel_conditions}
    Let $w: \mathbb{R}_{\geq 0} \to \mathbb{R}_{\geq 0}$ be a fixed function. The input kernel has the form of
    \begin{align} \label{tsne:input_kernel}
        \ki(x,x',\sigma) = \exp(-w(\sigma\norm{x-x'}^{\theta}))
    \end{align}
    Here, $\theta>0$ is a fixed constant. $w$ satisfies:

    1. $\forall t \geq 0$, $w'(t)>0$. $\lim_{t \to \infty} w(t) = \infty$.

    2. $w'(t)+tw''(t) \geq 0$ for all $t \geq 0$.

    3. The following integral converges:
    \begin{align}
        \int_0^{\infty} t^{d-1} w(t^{\theta}) \exp(-w(t^{\theta}))dt < \infty.
    \end{align}
\end{condition}
Since the all probability masses are normalized, we can assume without loss of generality that $w(0)=0$ from now on.
\begin{condition}[Output Kernel Conditions] \label{output_kernel_conditions}
    The output kernel depends only on the distance between two output data points $y$ and $y'$. In particular, $\ko(y,y') = k (\norm{y-y'})$ and:

    1. $k: \mathbb{R}_{\geq 0} \to \mathbb{R}^+$ is decreasing and bounded.

    2. $\ko$ satisfies
    \begin{align}
        \iint_{(0,\infty)^2} \ko(y,y') dy dy' < \infty.
    \end{align}

    3. $k$ has bounded derivative.

    4. $k'(0) = 0$.
\end{condition}

\subsection*{Setup and Notations} \label{tsne:sec:setup}
In this subsection, we provide an
introduction to the generalized t-SNE algorithm incorporating the new
terminology. Recall that the input of the generalized t-SNE consists of vectors
$X_1, \cdots, X_n \in \mathbb{R}^d$, with the purpose of finding vectors $Y_1,
    \cdots Y_n \in \mathbb{R}^s$ that serve as a good representation of the
original vectors. The algorithm starts by defining the following probability
masses:
\begin{align}
    p_{j|i} = \frac{\ki(X_i,X_j,\sigma_i)}{\sum_{k \neq i} \ki(X_i,X_k,\sigma_i)}.
\end{align}
Each $\sigma_i$ satisfies that for each $i = 1, \dots, n$,
\begin{align}
    -\sum_{j = 1, j \neq i}^n p_{j|i} \log p_{j|i} = \log \textbf{Perp} =: \log (n \rho).
\end{align}
$\textbf{Perp}$ is called the perplexity and $\rho \in (0,1)$ is a given constant. The purpose of $\sigma_i$ is to make sure that for each $i$, the probability distribution induced by $p_{j|i}$ has the same entropy across all $i$. Clearly, $p_{j|i}$ is not symmetric.
\begin{remark}
    For finite data set perplexity is usually set as a constant. Moreover, the scaling of perplexity with respect to $\log n$ is strictly necessary here. In \cite{murray2024largedatalimitsscaling}, it has been proved that if we let $\textbf{\normalfont Perp} = \log n^{\gamma} \rho$ for some constant parameter $\gamma$, we cannot prove convergence and the existence of equilibrium measures for neither $\gamma<1$ nor $\gamma >1$.
\end{remark}
We then define $p_{ii} = 0$ and for $1 \leq i \neq j \leq n$,
\begin{align}
    p_{ij} = \frac{p_{j|i} + p_{i|j}}{2n} = \frac{1}{2n}\left(\frac{\ki(X_i,X_j,\sigma_i)}{\sum_{k \neq i} \ki(X_i,X_k,\sigma_i)}+\frac{\ki(X_i,X_j,\sigma_j)}{\sum_{k \neq j} \ki(X_j,X_k,\sigma_j)}\right).
\end{align}
Obviously,
\begin{align}
    \sum_{\substack{j = 1 \\ j \neq i}}^n p_{j|i} = 1, \quad \sum_{1 \leq i \neq j \leq n} p_{ij} = 1.
\end{align}
As for the output vectors $Y_1, \cdots Y_n \in \mathbb{R}^s$, set for $1 \leq i \neq j \leq n$,
\begin{align}
    q_{ij} = \frac{\ko(Y_i,Y_j)}{\sum_{k \neq l}\ko(Y_k,Y_l)}.
\end{align}
The goal of the generalized t-SNE is to find $Y$ that minimizes the relative entropy
\begin{align}
    L_{n,\rho}(X,Y) = \sum_{1 \leq i \neq j \leq n} p_{ij} \log \left(\frac{p_{ij}}{q_{ij}}\right).
\end{align}
More precisely, the output of t-SNE $Y^*$ is defined as
\begin{align}
    Y^* = \argmin_{Y \in (\mathbb{R}^s)^n}L_{n,\rho}(X,Y).
\end{align}
In the expression of $L_n(X,Y)$, all $p_{ij}$ are given and $q_{ij}$ is a function of $Y_i$ and $Y_j$. Therefore, $Y^*$ can be found simply by using gradient descent. This completes the description of the algorithm itself. Now we introduce necessary notations for convergence and asymptotic analysis. Let $\mathcal{M}(\mathbb{R}^d)$ be the space of Borel probability measures on $\mathbb{R}^d$ and $\mu_X \in \mathcal{M}(\mathbb{R}^d)$ be a probability measure. Suppose all input $X = (X_1,\cdots,X_n)$ are i.i.d. random variables with probability measure $\mu_X$. Meanwhile, let $Y^* = (Y_1,\cdots,Y_n)$ be the output of the t-SNE algorithm. Consider the following empirical measure:
\begin{align}
    \mu_n := \frac{1}{n} \sum_{i=1}^n \delta(X_i,Y_i) \in \mathcal{M}(\mathbb{R}^d \times \mathbb{R}^s).
\end{align}
We are interested in its limiting measure as $n \to \infty$. To finish the setup, there are a few more remaining quantities we would like to introduce. Let $\mu \in \mathcal{M}(\mathbb{R}^d \times \mathbb{R}^s)$. For a given kernel $\ki(x,y)$, define $F_{\rho,\mu}(x,\sigma) : \mathbb{R}^d \times \mathbb{R}^+ \to \mathbb{R}$ as
\begin{align}
    F_{\rho,\mu}(x,\sigma) = \int \frac{\ki(x,x',\sigma)}{\int \ki(x,x',\sigma)d\mu(x')} \log \left(\frac{\ki(x,x',\sigma)}{\int \ki(x,x',\sigma)d\mu(x')}\right) d\mu(x') + \log \rho.
\end{align}
Note that the last $s$ variables in the integrands stay constant. It is not hard to see that $F_{\rho,\mu}(x,\sigma)$ is the continuous version of the entropy. Next, we extend the definition of $\sigma_i$ by defining $\sigma^*_{\rho,\mu} : \mathbb{R}^d \to \mathbb{R}$ the unique solution of the equation
\begin{align}
    F_{\rho,\mu}(x,\sigma) = 0.
\end{align}
The introduction of $\sigma^*$ aligns with the procedure of finding $\sigma_i$ that controls the entropy to be the same for all $i$. Of course, the existence and uniqueness of $\sigma^*$ is not trivial and will be discussed in section \ref{tsne:sec:existence_uniqueness}. For the symmetric probability mass, let $\psi: \mathbb{R}^d \to \mathbb{R}$ be a given function and $p_{\psi}: \mathbb{R}^d \times \mathbb{R}^d \to \mathbb{R}$ be defined by
\begin{align}
    p_{\psi}(x,x') = \frac{1}{2}\left( \frac{\ki(x,x',\psi(x))}{\int \ki(x,x',\psi(x))d\mu(x')} + \frac{\ki(x,x',\psi(x'))}{\int \ki(x,x',\psi(x'))d\mu(x)}\right).
\end{align}
For the extension on the output side, let $q: \mathbb{R}^s \times \mathbb{R}^s \to \mathbb{R}$ be given by
\begin{align}
    q(y,y') = \frac{\ko(y,y')}{\iint \ko(y,y')d\mu(y) d\mu(y')}.
\end{align}
Based on the construction of $p$ and $q$, define
\begin{align}
    I_{\rho}(\mu) = \iint p_{\sigma^*_{\rho,\mu}}(x,x') \log\left(\frac{p_{\sigma^*_{\rho,\mu}}(x,x')}{q(y,y')}\right) d\mu(x,y) d\mu(x',y').
\end{align}
Recall $\mu_X \in \mathcal{M}(\mathbb{R}^d)$ is an arbitrary given probability measure. We set $\mathcal{P}_X$ to be the set of probability measures on $\mathbb{R}^d \times \mathbb{R}^s$ whose marginal coincides with $\mu_X$. Formally speaking,
\begin{align}
    \mathcal{P}_X := \{\mu \in \mathcal{M}(\mathbb{R}^d \times \mathbb{R}^s): \text{the $\mathbb{R}^d$ marginal of $\mu$ is $\mu_X$}\}.
\end{align}
Lastly, define
\begin{align}
    \tilde{\mathcal{P}}_X := \left\{ \mu \in \mathcal{P}_X: \int (p_{\sigma^*_{\rho,\mu}}(x,x')-q(y,y'))\frac{\frac{\partial}{\partial y}\ko{(y,y')}}{\ko{(y,y')}} d\mu(x',y') = 0,\mu(x,y)\text{-a.e.}\right\}.
\end{align}
\section{Main Results} \label{tsne:sec:main_results}
Using the established setup, we are able to prove the same results as in
\cite{auffinger-tsne} with a broad class of input and output kernels. On top of
that, compared to the original paper, we no longer require the sample
probability measure $\mu_X$ to have $C^1$ density. We managed to relax the
condition to $C^0$.
\begin{theorem} \label{tsne:main_theorem_1}
    Assume $\mu_X$ has a continuous density $f(x)$ and compact support. Let $X_i \in \mathbb{R}^d$ be i.i.d random variables with probability measure $\mu_X$. If condition \ref{input_kernel_conditions} and condition \ref{output_kernel_conditions} hold, then for any $\rho \in (0,1)$,
    \begin{align}
        \lim_{n \to \infty} \inf_Y L_{n,\rho}(X,Y) = \inf_{\mu \in \tilde{\mathcal{P}}_X} I_{\rho}(\mu).
    \end{align}
    Moreover, there exist a measure $\mu^*$ and a sub-sequence $\{n_k\}$ such that $\mu_{n_k} \to \mu^*$ weakly and $I_{\rho}(\mu^*) = \inf_{\mu \in \tilde{\mathcal{P}}_X} I_{\rho}(\mu)$.
\end{theorem}
\begin{theorem} \label{tsne:main_theorem_2}
    Under the same assumptions of theorem \ref{tsne:main_theorem_1}, $\mu^*$ has compact support.
\end{theorem}
\begin{remark}
    In \cite{auffinger-tsne}, several lemmas are proved under the assumption that the sample distribution $\mu_X$ is sub-Gaussian. This can also be done for our work, if we replace condition \ref{input_kernel_conditions} with the following condition \ref{input_kernel_conditions_ext}:
    \begin{condition}[Alternative Input Kernel Conditions] \label{input_kernel_conditions_ext}
        Let $w: \mathbb{R}_{\geq 0} \to \mathbb{R}_{\geq 0}$ be some fixed function. The input kernel has the form of $\ki(x,x,\sigma) = \exp(-w(\sigma\norm{x-x'}^{\theta}))$. Here, $\sigma$ is used to control entropy (thus perplexity) and $\theta>0$ is a fixed constant. Define
        \begin{align} \label{tsne:kw_def}
            \mathcal{K}_w = \left\{\tilde{C} \in \mathbb{R}_{\geq 0} \bigg\rvert \exists \epsilon>0, \text{\normalfont s.t.} \lim_{n \to \infty} n^{-1/2} \exp \left(6 w\left(\left(\tilde{C} \sqrt{\log n} \right)^{\theta} \right) \right) = O(n^{-\epsilon}) \right\}.
        \end{align}
        Then $w$ satisfies:

        1. $\forall t \geq 0$, $w'(t)>0$. $\lim_{t \to \infty} w(t) = \infty$;

        2. $w'(t)+tw''(t) \geq 0$ for all $t \geq 0$;

        3. The following integral converges:
        \begin{align}
            \int_0^{\infty} t^{d-1} w(t^{\theta}) \exp(-w(t^{\theta}))dt < \infty;
        \end{align}

        4. $\sup \mathcal{K}_w > 0$, and $\forall t \geq 0, w'(t) \leq w(t)$.
    \end{condition}
    We will not include proofs for lemmas assuming condition \ref{input_kernel_conditions_ext} and sub-Gaussian samples, because such generalization will not carry onto the main theorems. One could adapt the proofs we present assuming $\mu_X$ being sub-Gaussian and condition \ref{input_kernel_conditions_ext}. However, it remains unknown that whether both theorem \ref{tsne:main_theorem_1} and \ref{tsne:main_theorem_2} hold assuming a more general data distribution.
\end{remark}
\section{The Generality of Conditions} \label{tsne:sec:generality_conditions}
Since the main goal of our work is to generalize the theoretical guarantees of
the t-SNE algorithm, it is important that we do not lose comprehensiveness when
assuming a concrete form of both the input and output kernels. Therefore, both
condition \ref{input_kernel_conditions} and \ref{output_kernel_conditions}
warrant an explanation.

For the formula of the input kernel in definition \ref{def:input_kernel}, we
start by realizing that $t$-SNE is about the relationship between all data
points. The insights we gain from the algorithm should be translation and
rotation invariant with respect to the coordinate system we choose. Hence, it
is only reasonable to assume that the input kernel is a function of
$\norm{x-x'}$ for each pair of input data points $x,x'$.

On the other hand, the weight we assign to each pair of data points is always
positive due to the nature of probability masses. We thus lose no generality by
assuming the input kernel has an exponential form because obviously, the
exponential function is a smooth bijection from $\mathbb{R}$ to $\mathbb{R}^+$.
The same above can be said for the output kernel.

Additionally, though it is coincidental that in the original version of the
algorithm $\sigma_i$ can be viewed as the variance of a Gaussian random
variable, in reality $\sigma_i$'s are only used to control the entropy of each
probability distribution induced by $x_i$. With that being said, we only need
to introduce $\sigma$ in such a way that it is able to control the entropy for
each $x$. As we can see in section \ref{tsne:sec:existence_uniqueness}, simply
multiplying $\sigma$ by $\norm{x-x'}^{\theta}$ is enough to achieve this goal.
Moreover, we have the convergence of $\sigma^*$ in section
\ref{tsne:sec:input_kernel_convergence}, providing the anchor of the following
results just as in \cite{auffinger-tsne}.

Combining all the above, we arrive at the form of the input kernel in
\eqref{tsne:input_kernel}, by adding a general function $w$ to the product
$\sigma \norm{x-x'}^{\theta}$. This gives us a kernel we can easily work on,
without losing any generality.

Besides the form of the input kernel, we also want to make sure that the
constraints we impose on the kernels in condition \ref{input_kernel_conditions}
and \ref{output_kernel_conditions} are also general enough for the significance
of our results. Since we frequently deal with integrals containing the input
and output kernels, their integrability is the bare minimum we should assume.
Along with integrability comes the requirement of their decay to 0 at infinity.
Apart from that, we do not add any excessive constraints to the kernels and
readers will see later in proofs that those assumptions are strictly necessary
for the results to hold.

Specifically, notice the second assumption of condition
\ref{input_kernel_conditions} is basically requiring the input kernel to decay
at least at a minimal speed by controlling its second order derivative. If we
let $w(t) = \log t$, we have $w'(t)+tw''(t) = 0$. Roughly speaking, this
implies that the main theorems hold as long as the input kernel decays at
polynomial rate or faster, which is a big step up than proving the convergence
for Gaussian kernels alone.
\section{Summary and Structure of Proofs} \label{tsne:sec:summary_structure}
Before we move onto mathematical details, we dedicate this short section to
give a summary of our work and the structure of the proofs. Since t-SNE assigns
a ``weight" to each pair of data points using Gaussian and $t$-distribution, we
propose the concept of generalized kernels so that the generalized t-SNE can
use a wide range of functions to define the weights, thus the subsequent
probability distributions associated with the input and output data points. A
natural question to ask is: under which conditions can we prove the convergence
of the generalized t-SNE algorithm towards an equilibrium measure? Our work
here answers the questions by providing a concrete formulation of the
generalized kernels as well as rigorous proofs of the convergence towards the
equilibrium measure.

The generality of the kernel formulation as well as the conditions we impose
are discussed in the previous section. For the proofs, we first establish the
existence and uniqueness of $\sigma^*$ as the asymptotic limit of $\sigma_i$'s
in section \ref{tsne:sec:existence_uniqueness}. Then we state the uniform
convergence of $\sigma^*$ in section \ref{tsne:sec:input_kernel_convergence}.
$\sigma^*$ serves as the anchor point of all subsequent convergence results.
With the existence and uniqueness of $\sigma^*$ as well as its convergence
established, we are able to prove the convergence towards the equilibrium
measure by proving inequalities from both sides, the details of which are
discussed in section \ref{tsne:sec:lower_bound} and \ref{tsne:sec:upper_bound}.
Finally the proofs are wrapped up in section \ref{tsne:sec:main_theorem_1} and
\ref{tsne:sec:main_theorem_2}, where we show that the equilibrium measure also
has compact support just as the input measure $\mu_X$.
\section{Existence and Uniqueness of \texorpdfstring{$\sigma^*$}{sigma-star}}\label{tsne:sec:existence_uniqueness}
In this section we establish the existence and uniqueness of $\sigma^*_{\rho,\mu}(x)$. Recall that $\sigma^*_{\rho,\mu}(x)$ is the unique solution of the equation
\begin{align}
    F_{\rho,\mu}(x,\sigma) = \int \frac{\ki(x,x',\sigma)}{\int \ki(x,x',\sigma)d\mu(x')} \log \left(\frac{\ki(x,x',\sigma)}{\int \ki(x,x',\sigma)d\mu(x')}\right) d\mu(x') + \log \rho = 0
\end{align}
Intuitively speaking, $\sigma^*$ is the limit function of $\sigma_i$'s that are used to control the entropy. The purpose of this section is to prove that such limit exists and is unique for each given $x$. The existence and uniqueness of $\sigma^*$ will become the foundation of everything we will prove afterwards. By condition \ref{input_kernel_conditions}, we can write
\begin{align} \label{big_F_formula}
      & F_{\rho,\mu}(x,\sigma) \nonumber                                                                                                                                                                                                                    \\
    = & \int \frac{\exp(-w(\sigma\norm{x-x'}^{\theta}))}{\int \exp(-w(\sigma\norm{x-x'}^{\theta})) d\mu(x')} \log\left(\frac{\exp(-w(\sigma\norm{x-x'}^{\theta}))}{\int \exp(-w(\sigma\norm{x-x'}^{\theta})) d\mu(x')}\right)d\mu(x') + \log \rho \nonumber \\
    = & -\int \frac{w(\sigma\norm{x-x'}^{\theta})\exp(-w(\sigma\norm{x-x'}^{\theta}))}{\int \exp(-w(\sigma\norm{x-x'}^{\theta})) d\mu(x')} d\mu(x') \nonumber                                                                                               \\
      & \qquad \qquad \qquad \qquad - \log \left(\int \exp(-w(\sigma\norm{x-x'}^{\theta}))d\mu(x')\right)  + \log \rho.
\end{align}
We start with the lemma that gives $F$ an increasing lower bound.
\begin{lemma} \label{existence_uniqueness:existence_lemma}
    If $\mu$ has a continuous, bounded density $f(x)$, then there exists a constant $C>0$ such that
    \begin{align}  \label{big_F_lowerbound_final}
        F_{\rho,\mu}(x,\sigma) \geq C \log \sigma - \log f(x) + \log \rho.
    \end{align}
\end{lemma}
\begin{proof}
    Notice that for any integrable function $h$ on $\mathbb{R}_{\geq 0}$, we have
    \begin{align}
         & \int_{\mathbb{R}^d} h(\norm{x}) dx = \int_{\mathbb{S}^{d-1}} \int_0^{\infty} r^{d-1}h(r) dr d\nu=S_{d-1} \int_0^{\infty} r^{d-1}h(r) dr,
    \end{align}
    where $\mathbb{S}^{d-1}$ is the $(d-1)$-dimensional sphere, $d\nu$ the corresponding surface measure, and $S_{d-1}$ the volume of the sphere. By Condition \ref{input_kernel_conditions}, we know that
    \begin{align}
        \int_0^{\infty} t^{d-1} w(t^{\theta}) \exp(-w(t^{\theta}))dt < \infty.
    \end{align}
    Since $w$ is increasing and $\lim_{t \to \infty} w(t) = \infty$, it is easy to see that
    \begin{align}
        \int_0^{\infty} t^{d-1} \exp(-w(t^{\theta}))dt < \infty.
    \end{align}
    Denote
    \begin{align}
        Z_d :=\int_{\mathbb{R}^d}\exp(-w(\norm{\tau}^{\theta})) d \tau = S_{d-1} \int_0^{\infty} t^{d-1}\exp(-w(t^{\theta})) dt.
    \end{align}
    Consider
    \begin{align}
          & \int \exp(-w(\sigma\norm{x-x'}^{\theta}))d\mu(x') = \int \exp(-w(\sigma\norm{x-x'}^{\theta})) f(x')dx' \nonumber \\
        = & Z_d \sigma^{-d/\theta}\int \frac{\exp(-w(\sigma\norm{x-x'}^{\theta}))}{Z_d \sigma^{-d/\theta}} f(x')dx'.
    \end{align}
    Because $\lim_{x \to \infty} w(x) = \infty$, the integrand of the above integral converges weakly to $\delta_x f(x')$. To prove this statement, notice that
    \begin{align}
               & \int \frac{\exp(-w(\sigma\norm{x-x'}^{\theta}))}{Z_d \sigma^{-d/\theta}} f(x')dx' \nonumber                                          \\
        =      & \frac{1}{Z_d} \int \exp\left(-w \left(\norm{\sigma^{1/\theta}(x-x')}^{\theta}\right)\right) f(x') d\sigma^{1/\theta}(x-x') \nonumber \\
        \equiv & \frac{1}{Z_d} \int \exp(-w(\norm{\tau}^{\theta})) f(x-\tau \sigma^{-1/\theta}) d\tau.
    \end{align}
    Since $f$ is bounded, $\exp(-w(\norm{\tau}^{\theta})) \sup f$ is integrable. By Dominant Convergence Theorem and the continuity of $f$, we have
    \begin{align}
          & \lim_{\sigma \to \infty} \frac{1}{Z_d} \int \exp(-w(\norm{\tau}^{\theta})) f(x-\tau \sigma^{-1/\theta}) d\tau \nonumber \\
        = & \frac{1}{Z_d} \int \lim_{\sigma \to \infty} \exp(-w(\norm{\tau}^{\theta})) f(x-\tau \sigma^{-1/\theta}) d\tau \nonumber \\
        = & f(x) \frac{1}{Z_d} \int \exp(-w(\norm{\tau}^{\theta})) d\tau \nonumber                                                  \\
        = & f(x).
    \end{align}
    There exists some $\epsilon(\sigma) \to 0$ as $\sigma\to \infty$ and a constant $C>0$, that
    \begin{align} \label{big_f_lowerbound_1}
        -\log \left(\int \exp(-w(\sigma\norm{x-x'}^{\theta}))d\mu(x')\right) = & -\log\left(Z_d \sigma^{-d/\theta}\right) - \log f(x)+ \epsilon(\sigma) \nonumber \\
        \geq                                                                   & C\log \sigma - \log f(x).
    \end{align}
    Next, since $f$ is bounded, there exists $M>0$ such that $f \leq M$.
    \begin{align} \label{big_F_lowerbound_2}
             & \int \frac{w(\sigma\norm{x-x'}^{\theta})\exp(-w(\sigma\norm{x-x'}^{\theta}))}{\int \exp(-w(\sigma\norm{x-x'}^{\theta})) d\mu(x')} d\mu(x') \nonumber                                                                  \\
        \leq & M \frac{\int w(\sigma\norm{x-x'}^{\theta})\exp(-w(\sigma\norm{x-x'}^{\theta}))dx'}{\int \exp(-w(\sigma\norm{x-x'}^{\theta}))f(x') dx'} \nonumber                                                                      \\
        =    & M \frac{\sigma^{-d/\theta}\int w(\norm{\sigma^{1/\theta}(x-x')}^{\theta})\exp(-w(\norm{\sigma^{1/\theta}(x-x')}^{\theta}))d\sigma^{1/\theta}(x-x')}{ \int \exp(-w(\sigma \norm{(x-x')}^{\theta}))f(x') dx'} \nonumber \\
        =    & M \frac{\sigma^{-d/\theta}\int w(\norm{\tau}^{\theta})\exp(-w(\norm{\tau}^{\theta}))d\tau}{ \int \exp(-w(\sigma \norm{(x-x')}^{\theta}))f(x') dx'} \nonumber                                                          \\
        % =    & M \frac{\int w(\norm{t}^{\theta})\exp(-w(\norm{t}^{\theta}))dt}{Z_d \int \exp(-w(\sigma \norm{(x-x')}^{\theta}))f(x') dx'} \nonumber                                                                                  \\
        \to  & \frac{M}{Z_d f(x)} \int w(\norm{\tau}^{\theta})\exp(-w(\norm{\tau}^{\theta}))d\tau < \infty.
    \end{align}
    The numerator converges because of condition \ref{input_kernel_conditions}. As the bound in (\ref{big_F_lowerbound_2}) does not depend on $\sigma$, we can combine (\ref{big_F_formula}), (\ref{big_f_lowerbound_1}) and (\ref{big_F_lowerbound_2}) and deduce that
    \begin{align}
        F_{\rho,\mu}(x,\sigma) \geq C \log \sigma - \log f(x) + \log \rho.
    \end{align}
\end{proof}
This lemma proves the existence of $\sigma^*_{\rho,\mu}(x)$. Indeed, by \eqref{big_F_formula}, apparently
\begin{align} \label{existence_uniqueness:limit_1}
    \lim_{\sigma \to 0} F_{\rho,\mu}(x,\sigma) = \log \rho < 0
\end{align}
since $\rho \in (0,1)$. Now with lemma \ref{existence_uniqueness:existence_lemma}, we get
\begin{align} \label{existence_uniqueness:limit_2}
    \lim_{\sigma \to \infty} F_{\rho,\mu}(x,\sigma) = \infty.
\end{align}
By continuity of $f$ and intermediate value theorem, $\sigma^*_{\rho,\mu}(x)$ exists. To prove the uniqueness, we are left to show $F_{\rho,\mu}$ is strictly increasing. The next lemma will do just that.
\begin{lemma} \label{existence_uniqueness:uniqueness_lemma}
    Given the same condition as in lemma \ref{existence_uniqueness:existence_lemma}, $F_{\rho,\mu}(x,\sigma)$ is strictly increasing in $\sigma$.
\end{lemma}
To tackle this proof, we need one additional lemma inspired by theorem 7 in \cite{liu2000inequ}.
\begin{lemma} \label{miracle_lemma_ext}
    Let $f(x),g(x),h(x)$ be continuous functions $\mathbb{R}^d \to \mathbb{R}^+$ such that for all $x,y \in \mathbb{R}^d$
    \begin{align} \label{miracle_condition_ext}
        \left(g(x)-g(y)\right)\left(\frac{f(y)}{h(y)}-\frac{f(x)}{h(x)}\right) \geq 0.
    \end{align}
    Let $\mu$ be an arbitrary Lebesgue measure and assume all integrals in this lemma converge. Then the inequality
    \begin{align} \label{miracle_inequ_ext}
        \frac{\int_{\mathbb{R}^d} f(x)d\mu}{\int_{\mathbb{R}^d} h(x)d\mu} \geq \frac{\int_{\mathbb{R}^d} f(x)g(x)d\mu}{\int_{\mathbb{R}^d} h(x)g(x)d\mu}
    \end{align}
    holds. If (\ref{miracle_condition_ext}) reverses, then (\ref{miracle_inequ_ext}) reverses.
\end{lemma}
\begin{proof}
    The proof of this lemma is similar to the one in \cite{liu2000inequ}. It suffices to show that
    \begin{align}
        \int_{\mathbb{R}^d} f(x)d\mu \int_{\mathbb{R}^d} h(x)g(x)d\mu \geq \int_{\mathbb{R}^d} h(x)d\mu \int_{\mathbb{R}^d} f(x)g(x)d\mu.
    \end{align}
    The above inequality is equivalent to
    \begin{align}
        \int_{\mathbb{R}^d} f(x)d\mu(x) \int_{\mathbb{R}^d} h(y)g(y)d\mu(y) \geq \int_{\mathbb{R}^d} h(x)d\mu(x) \int_{\mathbb{R}^d} f(y)g(y)d\mu(y).
    \end{align}
    That means we need to show
    \begin{align}
        \mathcal{G}:=\int_{\mathbb{R}^d} \int_{\mathbb{R}^d} g(y)(f(x)h(y)-f(y)h(x))d\mu(x)d\mu(y) \geq 0.
    \end{align}
    Swap $x$ and $y$ in the formula of $\mathcal{G}$ and we have an equivalent expression:
    \begin{align}
        \mathcal{G}=\int_{\mathbb{R}^d} \int_{\mathbb{R}^d} g(x)(f(y)h(x)-f(x)h(y))d\mu(x)d\mu(y).
    \end{align}
    Adding the last two formulas together yields
    \begin{align}
        2\mathcal{G} = & \int_{\mathbb{R}^d} \int_{\mathbb{R}^d} (g(x)-g(y))(f(y)h(x)-f(x)h(y))d\mu(x)d\mu(y) \nonumber                                              \\
        =              & \int_{\mathbb{R}^d} \int_{\mathbb{R}^d} h(x)h(y)\left(g(x)-g(y)\right)\left(\frac{f(y)}{h(y)}-\frac{f(x)}{h(x)}\right)d\mu(x)d\mu(y) \geq 0
    \end{align}
    by the hypothesis of the lemma. This finishes the proof.
\end{proof}
\begin{proof}[Proof of lemma \ref{existence_uniqueness:uniqueness_lemma}]
    We compute the derivative of $F_{\rho,\mu}(x,\sigma)$ directly using \eqref{big_F_formula}. The following simplified
    notations will be used in the first several steps:
    \begin{align}
        w = w(\sigma\norm{x-x'}^{\theta}), \quad w'=w'(\sigma\norm{x-x'}^{\theta}), \quad \exp(\cdot) = \exp(-w(\sigma\norm{x-x'}^{\theta})).
    \end{align}
    Omitting $d\mu(x')$ in the integration for cleaner writing, and we have
    \begin{align}
          & \frac{\partial}{\partial \sigma} F_{\rho,\mu}(x,\sigma) \nonumber                                                                                                                                    \\
        = & -\int \frac{w'\norm{x-x'}^{\theta}\exp(\cdot)}{\int \exp(\cdot)} -\int \frac{w\exp(\cdot)'}{\int \exp(\cdot)} +
        \int \frac{w \exp(\cdot) \left(\int \exp(\cdot)\right)'}{\left(\int \exp(\cdot)\right)^2} - \frac{\left(\int \exp(\cdot)\right)'}{\int \exp(\cdot)} \nonumber                                            \\
        = & -\int \frac{w\exp(\cdot)'}{\int \exp(\cdot)} +
        \int \frac{w \exp(\cdot) \left(\int \exp(\cdot)\right)'}{\left(\int \exp(\cdot)\right)^2} \quad \text{\normalfont (The first and fourth terms cancel.)} \nonumber                                        \\
        = & \int \frac{w\cdot w' \cdot \norm{x-x'}^{\theta} \cdot \exp(\cdot)}{\int \exp(\cdot)}-\frac{\int w \cdot \exp(\cdot) \int w'\cdot \norm{x-x'}^{\theta} \exp(\cdot)}{\left(\int \exp(\cdot)\right)^2}.
    \end{align}
    We are able to move differentiation inside the integral because of dominant convergence theorem. To show that $\frac{\partial}{\partial \sigma} F_{\rho,\mu}(x,\sigma) \geq 0$, it suffices to show
    \begin{align} \label{monotone_inequ}
        \int w\cdot w' \cdot \norm{x-x'}^{\theta} \cdot \exp(\cdot) \int \exp(\cdot) \geq \int w \cdot \exp(\cdot) \int w'\cdot \norm{x-x'}^{\theta} \exp(\cdot).
    \end{align}
    For each $x,y,x' \in \mathbb{R}$ if the equality
    \begin{align} \label{monotone_condition}
         & (w'(\sigma\norm{x-x'}^{\theta})\norm{x-x'}^{\theta}-w'(\sigma\norm{y-x'}^{\theta})\norm{y-x'}^{\theta}) \nonumber                                   \\
         & \qquad \qquad \qquad \qquad \qquad \qquad \cdot \left(\frac{1}{w(\sigma\norm{y-x'}^{\theta})}-\frac{1}{w(\sigma\norm{x-x'}^{\theta})}\right) \geq 0
    \end{align}
    is true, then setting $f(x) = \exp(\cdot)$, $h(x) = w \cdot \exp(\cdot)$ and $g(x) = w'\cdot \norm{x-x'}^{\theta}$ in lemma \ref{miracle_lemma_ext} shows that (\ref{monotone_inequ}) holds. Let $\norm{x-x'}^{\theta} > \norm{y-x'}^{\theta}$ without loss of generality. Then by monotonicity of $w$, we have $w(\sigma\norm{x-x'}^{\theta})>w(\sigma\norm{y-x'}^{\theta})$. For (\ref{monotone_condition}) to be true, we are left to show
    \begin{align}
        w'(\sigma\norm{x-x'}^{\theta})\norm{x-x'}^{\theta}-w'(\sigma\norm{y-x'}^{\theta})\norm{y-x'}^{\theta} \geq 0.
    \end{align}
    Since $\sigma >0$, the above inequality is equivalent to:
    \begin{align} \label{generalized_condition_inequ}
        w'(\sigma\norm{x-x'}^{\theta})\sigma \norm{x-x'}^{\theta}-w'(\sigma\norm{y-x'}^{\theta}) \sigma \norm{y-x'}^{\theta} \geq 0.
    \end{align}
    Let $\tilde{w}(x) = x w'(x)$. Based on the second assumption in Condition \ref{input_kernel_conditions} and mean value theorem, $\tilde{w}$ satisfies for all $a,b\geq 0$, $a>b$, there exists $c \in [b,a]$ that
    \begin{align}
        \frac{aw'(a) - bw'(b)}{a-b} = w'(c) + cw''(c) \geq 0.
    \end{align}
    This proves \eqref{generalized_condition_inequ} by letting $a = \sigma\norm{x-x'}^{\theta}$ and $b = \sigma\norm{y-x'}^{\theta}$, as desired.
\end{proof}
We can now formally state and prove the main result of this section.
\begin{proposition}
    Let $\mu \in \mathcal{M}(\mathbb{R}^d)$ with $C^0$ density $f$ such that $f$ is bounded. The equation $F_{\rho,\mu}(x,\sigma) = 0$ has a unique solution $\sigma^*_{\rho,\mu}(x)$ for all $\rho \in (0,1)$ and $x \in \mathbb{R}^d$. Moreover, $x \mapsto \sigma^*_{\rho,\mu}(x)$ is smooth.
\end{proposition}
\begin{proof}
    The existence and uniqueness can be proven by combining \eqref{existence_uniqueness:limit_1}, \eqref{existence_uniqueness:limit_2}, lemma \ref{existence_uniqueness:existence_lemma} and lemma \ref{existence_uniqueness:uniqueness_lemma}. The smoothness of $\sigma^*_{\rho,\mu}(x)$ follows from the implicit function theorem.
\end{proof}

\section{Proofs of the Main Theorems} \label{tsne:sec:proofs}
From now on, we will assume both condition \ref{input_kernel_conditions} and
condition \ref{output_kernel_conditions} throughout this section. We will also
assume that the distribution $\mu_X$ that we draw samples from has $C^0$
density with bounded density and gradient of density. The assumptions on the
sample distribution $\mu_X$ are inherited directly from section
\ref{tsne:sec:existence_uniqueness} to avoid lengthy theorem and lemma
statements.
\subsection{Convergence Lemmas for Input Kernels} \label{tsne:sec:input_kernel_convergence}
The main difficulties of proving the convergence lemmas related to the input
kernel is the existence and uniqueness of $\sigma^*_{\rho,\mu}(x)$, which we
have already established in section \ref{tsne:sec:existence_uniqueness}. The
rest of lemmas are similar, if not easier than the corresponding results in
\cite{auffinger-tsne}. We will list the important lemmas here without proofs.
The reader can easily adapt the proofs in \cite{auffinger-tsne} to our case.

First we have the two following important propositions about the uniform
convergence of $F_{\mu_n}(x,\sigma)$ and $\sigma_{\mu_n}^*(x)$.
\begin{proposition} \label{tsne:uniform_convergence_prop}
    Let $\mu_X$ be a probability measure with compact support, and $X_1,\dots,X_n$ be i.i.d. random variables drawn from $\mu$. $\mu_n = (1/n)\sum_i \delta_{X_i}$ is the empirical measure. Then there exists a constant $\epsilon>0$, that almost surely in $X_1,\dots,X_n$,
    \begin{align}
        \sup_{\substack{\sigma^{-1} \in [\delta^{\theta},n^{\theta/2}] \\ i = 1,\dots, n}} \left|F_{\mu_n}(X_i,\sigma)-F_{\mu}(X_i,\sigma)\right| = O(n^{-\epsilon}).
    \end{align}
\end{proposition}
\begin{proposition} \label{tsne:uniform_convergence_prop_2}
    Suppose $\mu_X$ has compact support and $C^0$ density. Then almost surely, as $n \to \infty$,
    \begin{align}
        \sup_{i=1,\dots,n} \left|\sigma_{\mu_n}^*(X_i)-\sigma_{\mu}^*(X_i)\right| \to 0.
    \end{align}
    In particular,
    \begin{align}
        \left|\sigma_{\mu_n}^*(X_i)-\sigma_{\mu}^*(X_i)\right| = O\left(\left|F_{\mu_n}(X_i,\sigma)-F_{\mu}(X_i,\sigma)\right|\right).
    \end{align}
\end{proposition}
Next two lemmas can be borrowed directly from \cite{auffinger-tsne} as they irrelevant to the generalized kernels.
\begin{lemma} \label{mu_star_lemma}
    There exits $\mu^*$ such that
    \begin{align}
        I_{\rho}(\mu^*) = \inf_{\mu \in \tilde{\mathcal{P}}_X} I_{\rho}(\mu)
    \end{align}
\end{lemma}
To further simplify up the notations, define
\begin{align} \label{tsne:convergence_lemmas:notation}
    \sigma_n(x) = & \sigma_{\rho,\mu_n}^*(x), \sigma(x) = \sigma_{\rho,\mu_X}^*(x), \nonumber                                          \\
    f(x,x') =     & f_{\sigma}(x,x') =: \exp(-w(\sigma\norm{x-x'}^{\theta})), f_n(x,x') =  f_{\sigma_n}(x,x'), \nonumber               \\
    p'_n(x,x') =  & \frac{f_n(x,x')}{\int f_n(x,x')d\mu_n(x') - 1/n}, \quad p'(x,x') = \frac{f(x,x')}{\int f(x,x')d\mu(x')}, \nonumber \\
    p_n(x,x') =   & \frac{1}{2}(p'_n(x,x')+p'_n(x',x)), \quad p(x,x') = \frac{1}{2}(p'(x,x')+p'(x',x)).
\end{align}
Note that $p_{j|i} = p'_n(X_i,X_j)/n$ and $p_{ij} = p_n(X_i,X_j)/n^2$, and for $\sigma^{-1} \in [\delta^{\theta},n^{\theta/2}]$ we have
\begin{align} \label{tsne:important_bounds}
    \int f(x,x')d\mu(x'), \quad \int f_n(x,x')d\mu_n(x') \geq M
\end{align}
for some constant $M>0$.

The next few lemmas cover the rest of convergence results for the input
kernels. The proofs can be easily adapted from \cite{auffinger-tsne}. We omit them here for
brevity.
\begin{lemma} \label{tsne:integrand_lemma_2}
    Let $\mu_X$ have compact support. There exists $\epsilon>0$ such that as $n \to \infty$,
    \begin{align}
        \sup_{i,j} \left|p'_n(X_i,X_j) - p'(X_i,X_j)\right| = & O(n^{-\epsilon}), \nonumber \\
        \sup_{i,j} \left|p_n(X_i,X_j) - p(X_i,X_j)\right| =   & O(n^{-\epsilon}) \nonumber.
    \end{align}
\end{lemma}
\begin{lemma} \label{entropy_lemma_1}
    As $n \to \infty$, we have
    \begin{align}
        \sum_{i \neq j} p_{ij} \log p_{ij} = \iint p(x,x')\log p(x,x')d\mu d\mu - 2 \log n + o(1).
    \end{align}
\end{lemma}
\subsection{Lower Bound for \texorpdfstring{$\inf_{Y \in (\mathbb{R}^s)^n}L_{n,\rho}(X,Y)$}{L(X,Y)}} \label{tsne:sec:lower_bound}
We write $\ko(y,y') = g(y,y') = k(\norm{y-y'})$ from now on for simpler notations.
\begin{lemma} \label{lower_bound_lemma_1}
    There exists a set of points $y^*(X)$ such that
    \begin{align}
        \inf_{Y \in (\mathbb{R}^s)^n}L_{n,\rho}(X,Y) = L_{n,\rho}(X,y^*).
    \end{align}
\end{lemma}
\begin{proof}
    Let $k \in \{1,\cdots, n\}$. We are going to show that
    \begin{align}
        \lim_{|y_k| \to \infty} L_{n,\rho}(X,y) = \infty.
    \end{align}
    By definition, $L$ is the relative entropy and
    \begin{align}
        L_{n,\rho}(X,y) = \sum_{i \neq j} p_{ij}\log \left(\frac{p_{ij}}{q_{ij}}\right) = \sum_{i \neq j} p_{ij}\log p_{ij} - \sum_{i \neq j} p_{ij}\log q_{ij}.
    \end{align}
    For $n$ given and fixed, it suffices to show that
    \begin{align}
        \sum_{i \neq j} p_{ij}\log (q_{ij}) = \sum_{i \neq j} p_{ij}\log \left(\frac{g(y_i,y_j)}{\sum_{m \neq l}g(y_m,y_l)}\right) \to -\infty.
    \end{align}
    Recall that $g(y,y')$ is bounded by condition \ref{output_kernel_conditions}. We have
    \begin{align}
             & -\sum_{i \neq j} p_{ij} \log \left(\sum_{m \neq l}g(y_m,y_l) \right) = -\log \left(\sum_{m \neq l}g(y_m,y_l) \right) \nonumber \\
        \leq & -\sum_{i \neq j} p_{ij} \log \left(\sum_{m\neq k, l \neq k, m \neq l}g(y_m,y_l) \right) < \infty.
    \end{align}
    On the other hand, we have $g(y_i,y_k) = k (\norm{y_i-y_k}) \to 0$ as $|y_k| \to \infty$ because $\norm{y_i-y_k} \to \infty$ and $k$ satisfies
    \begin{align}
        \iint_{(0,\infty)^2} k(\norm{x-x'}) dxdx' < \infty.
    \end{align}
    This means
    \begin{align}
        \lim_{|y_k| \to \infty}\sum_{i \neq j} p_{ij} \log (g(y_i,y_j)) \to -\infty.
    \end{align}
    Hence, the minimum of the $L_{n,\rho}(X,y)$ is achieved on a bounded set, the lemma follows by the continuity of $L_{n,\rho}(X,y)$.
\end{proof}
Before we move on, we would like to take a small deviation and compute the
partial derivative of the relative entropy. This is important because it will
impact the property of the set $\tilde{\mathcal{P}}_X$.
\begin{lemma} \label{derivative_lemma}
    Let $X \in (\mathbb{R}^d)^n$ and $Y = (y_1,\dots,y_n)\in (\mathbb{R}^s)^n$ be the input and output of the t-SNE algorithm, then we have
    \begin{align}
        \frac{\partial L_{n, \rho}(X,Y)}{\partial y_i} = -\sum_{j=1, j \neq i}^n (p_{ij}-q_{ij}) \frac{\frac{\partial}{\partial y_i}g(y_i,y_j)}{g(y_i,y_j)}.
    \end{align}
\end{lemma}
\begin{proof}
    By definition,
    \begin{align}
          & L_{n, \rho}(X,Y) = \sum_{i \neq j} p_{ij} \log \left(\frac{p_{ij}}{q_ij}\right) \nonumber \\
        = & \sum_{i \neq j} p_{ij} \log p_{ij} - \sum_{i \neq j} p_{ij} \log q_{ij}.
    \end{align}
    The first summation does not depend on $y_i$ and will vanish when taking the partial derivative. Moreover, we can further breakdown the second summation:
    \begin{align}
          & - \sum_{i \neq j} p_{ij} \log q_{ij} = - \sum_{i \neq j} p_{ij} \log g(y_i,y_j) + \sum_{i \neq j} p_{ij} \log \left(\sum_{k \neq l} g(y_k,y_l) \right)\nonumber \\
        = & - \sum_{i \neq j} p_{ij} \log g(y_i,y_j) + \log \left(\sum_{k \neq l} g(y_k,y_l)\right).
    \end{align}
    Taking derivative of the last expression and note that only terms that contain $y_i$ remain. This means
    \begin{align} \label{L_derivative}
        \frac{\partial L_{n, \rho}(X,Y)}{\partial y_i} = & \sum_{j \neq i, j = 1}^n \left(-p_{ij}\frac{\frac{\partial}{\partial y_i}g(y_i,y_j)}{g(y_i,y_j)} + \frac{1}{\sum_{k \neq l} g(y_k,y_l)} \frac{\partial}{\partial y_i} g(y_i,y_j)\right) \nonumber                            \\
        =                                                & \sum_{j \neq i, j = 1}^n \left(-p_{ij}\frac{\frac{\partial}{\partial y_i}g(y_i,y_j)}{g(y_i,y_j)} + \frac{g(y_i,y_j)}{\sum_{k \neq l} g(y_k,y_l)} \frac{\frac{\partial}{\partial y_i}g(y_i,y_j)}{g(y_i,y_j)}\right) \nonumber \\
        =                                                & -\sum_{j=1, j \neq i}^n (p_{ij}-q_{ij}) \frac{\frac{\partial}{\partial y_i}g(y_i,y_j)}{g(y_i,y_j)}.
    \end{align}
\end{proof}
We can easily do a sanity test by letting $g(y_i,y_j) = (1+\norm{y_i-y_j}^2)^{-1}$ and get formula (A.1) in \cite{auffinger-tsne}. The derivative in lemma \ref{derivative_lemma} explains why
$\tilde{\mathcal{{P}}}_X$ has that specific form.
\begin{lemma} \label{lower_bound_lemma_2}
    Assume $\mu_X$ has compact support. The following holds:
    \begin{align}
        \liminf_{n \to \infty}\sum_{i \neq j} p_{ij} \log \left(\frac{p_{ij}}{q_{ij}}\right) \geq \iint p(x,x')\log \frac{p(x,x')}{q(y,y')} d\mu d\mu.
    \end{align}
\end{lemma}
\begin{proof}
    Just as in lemma \ref{derivative_lemma}, we first break down the summation:
    \begin{align}
        \sum_{i \neq j} p_{ij} \log \left(\frac{p_{ij}}{q_{ij}}\right) = \sum_{i \neq j} p_{ij} \log p_{ij} - \sum_{i \neq j} p_{ij} \log q_{ij}.
    \end{align}
    The convergence of the first summation has been established in lemma \ref{entropy_lemma_1}. We just need to deal with the second summation. Note that we now have
    \begin{align}
        \mu_n = \frac{1}{n} \sum \delta_{X_i} \times \delta_{y_i}
    \end{align}
    if $y_i$ is involved. By the definition of $q_{ij}$, we have
    \begin{align} \label{entropy_esti_2}
          & -\sum_{i \neq j}p_{ij} \log(q_{ij}) \nonumber                                                                                             \\
        = & \frac{1}{n^2}\sum_{i \neq j} p_n(X_i,X_j) \log g(y_i,y_j) + \log \left(\frac{1}{n^2}\sum_{i \neq j} g(y_i,y_j)\right) +2 \log n \nonumber \\
        = & \iint p_n(x,x')g(y,y')d\mu_n d\mu_n + \log \left(\iint g(y,y')d\mu_n d\mu_n - \frac{1}{n}\right)+2\log n.
    \end{align}
    By Fatou's lemma and lemma \ref{tsne:integrand_lemma_2}, we have
    \begin{align} \label{fatou_esti_1}
        \liminf_{n \to \infty} p_n(x,x')\log g(y,y') d\mu_n d\mu_n \geq  \iint p(x,x')\log g(y,y') d\mu d\mu.
    \end{align}
    Similarly, using Fatou lemma again and we have
    \begin{align} \label{fatou_esti_2}
        \liminf_{n \to \infty} \iint g(y,y') d\mu_n d\mu_n - \frac{1}{n} \geq \iint g(y,y')d\mu d\mu.
    \end{align}
    The $1/n$ is the duplicated diagonal term in the empirical measure and will obviously vanish as we take $n \to \infty$. Moreover, because of condition \ref{output_kernel_conditions}, the integrals on the RHS of the last two inequalities converge. In addition, since $\log$ is a continuous and increasing function, (\ref{fatou_esti_2}) implies
    \begin{align} \label{fatou_esti_3}
        \liminf_{n \to \infty} \log \left(\iint g(y,y') d\mu_n d\mu_n - \frac{1}{n}\right) \geq \log \left(\iint g(y,y')d\mu d\mu \right).
    \end{align}
    The proof is finished by combining (\ref{entropy_esti_2}), (\ref{fatou_esti_1}), (\ref{fatou_esti_3}) and lemma \ref{entropy_lemma_1}. The $2 \log n$ cancelled out because of lemma \ref{entropy_lemma_1}.
\end{proof}
\begin{proposition} \label{lower_bound_prop}
    There exists a sub-sequence $n(k)$ such that
    \begin{align}
        \liminf_{k \to \infty} \inf_{Y \in (\mathbb{R}^s)^{n(k)}} L_{n(k),\rho}(X,Y) \geq I_{\rho}(\mu).
    \end{align}
\end{proposition}
\begin{proof}
    Let $y^*$ be given in the lemma and define
    \begin{align}
        \mu_n = \frac{1}{n} \sum \delta_{X_i} \times \delta_{y_i}.
    \end{align}
    We can easily show that $\mu_n$ converges weakly to a measure $\mu$. By lemma \ref{appendix_lemma} in section \ref{tsne:sec:extra_proofs}, we know that the limiting measure $\mu \in \tilde{\mathcal{P}}_X$. The proof is complete by using lemma \ref{lower_bound_lemma_1} and lemma \ref{lower_bound_lemma_2}.
\end{proof}
\subsection{Upper Bound for \texorpdfstring{$\inf_{Y \in (\mathbb{R}^s)^n}L_{n,\rho}(X,Y)$}{L(X,Y)}} \label{tsne:sec:upper_bound}
For convenience, define
\begin{align}
    q_n(y,y') = \frac{g(y,y')}{\iint g(y,y')d \mu_n(y) d \mu_n (y') - \frac{1}{n}k(0)}.
\end{align}
Here $k(0) = g(y,y)$ is the value of the output kernel when both inputs are equal. Since multiplying a constant to $g$ does not change anything, we can just treat $k(0)=1$ without loss of generality as everything is normalized.
By lemma \ref{mu_star_lemma} there exists $\mu^*$ such that $I(\mu^*) = \inf_{\mu}I_{\rho}(\mu)$. Let $(X_1,Y_1), \dots, (X_n,Y_n)$ be i.i.d. random variables drawn from $\\mu^*$ and $\mu_n = \frac{1}{n} \sum_{i=1}^n \delta_{(X_i,Y_i)}$ their empirical measure. $\mu_n$ converges to $\mu^*$ weakly.
\begin{lemma} \label{int_approx_lemma_2}
    As $n \to \infty$, we have
    \begin{align}
        \sum_{i \neq j} \log \left(\frac{p_{ij}}{q_{ij}}\right) = \iint p_n(x,x') \log \left(\frac{p_n(x,x')}{q_n(y,y')}\right) d\mu_n d\mu_n + o(1).
    \end{align}
\end{lemma}
\begin{proof}
    We begin with the familiar computation
    \begin{align}
        \sum_{i \neq j} p_{ij} \log \left(\frac{p_{ij}}{q_{ij}}\right) = & \frac{1}{n^2}\sum_{i \neq j}p_n(X_i,X_j) \log \left(\frac{p_n(X_i,X_j)}{q_n(X_i,X_j)}\right) \nonumber          \\
        =                                                                & \iint p_n(x,x') \log \left(\frac{p_n(X_i,X_j)}{q_n(X_i,X_j)}\right) d \mu_n d \mu_n + \text{\normalfont Err}_n,
    \end{align}
    where
    \begin{align}
          & \text{\normalfont Err}_n \nonumber                                                                                                                                                                                                   \\
        = & -\frac{1}{n^2}\sum_{i=1}^n\frac{1}{\frac{1}{n}\sum_{k \neq i}f_n(X_i,X_k)}\left(\log \left(\frac{1}{\frac{1}{n}\sum_{k \neq i}f_n(X_i,X_k)}\right)-\log \left(\frac{1}{\frac{1}{n}\sum_{k \neq l}g(Y_k,Y_l)}\right)\right) \nonumber \\
        = & -\frac{1}{n^2}\sum_{i=1}^n\frac{1}{\frac{1}{n}\sum_{k \neq i}f_n(X_i,X_k)}\log \left(\frac{1}{\frac{1}{n}\sum_{k \neq i}f_n(X_i,X_k)}\right) \nonumber                                                                               \\
          & \qquad \qquad \qquad + \frac{1}{n^2}\sum_{i=1}^n\frac{1}{\frac{1}{n}\sum_{k \neq i}f_n(X_i,X_k)}\log \left(\frac{1}{\frac{1}{n}\sum_{k \neq l}g(Y_k,Y_l)}\right).
    \end{align}
    By bounds in \eqref{tsne:important_bounds}, it is easy to see that the first summation in the error term is $o(1)$ and we can write
    \begin{align}
        \text{\normalfont Err}_n = o(1) + \frac{1}{n^2}\sum_{i=1}^n\frac{1}{\frac{1}{n}\sum_{k \neq i}f_n(X_i,X_k)}\log \left(\frac{1}{\frac{1}{n}\sum_{k \neq l}g(Y_k,Y_l)}\right).
    \end{align}
    For the second term, it is already clear that
    \begin{align}
        \frac{1}{n^2}\sum_{i=1}^n\frac{1}{\frac{1}{n}\sum_{k \neq i}f_n(X_i,X_k)} = o(1).
    \end{align}
    On the other hand, we have by Fatou's lemma:
    \begin{align}
        \liminf_{n \to \infty}\frac{1}{n^2} \sum_{k \neq l}g(Y_k,Y_l) = \liminf_{n \to \infty} \iint g(y,y') d\mu_n d\mu_n - \frac{1}{n} \geq \iint g(y,y') d\mu^* d\mu^* < \infty,
    \end{align}
    by condition \ref{output_kernel_conditions}. We can see that the second term in $\text{\normalfont Err}_n$ is also $o(1)$, which immediately implies $\text{\normalfont Err}_n = o(1)$ and completes the proof.
\end{proof}
The next lemma is needed to close up the lower bound part of the main theorem.
\begin{lemma} \label{output_kernel_int_lemma}
    As $n \to \infty$, we have
    \begin{align}
        \iint g(y,y')d\mu_n(y) d\mu_n(y) - \frac{1}{n} = \iint g(y,y')d\mu^* d\mu^* + o(1).
    \end{align}
\end{lemma}
\begin{proof}
    Since $\mu_n \to \mu$ weakly, it is tight by Prokhorov's theorem. For arbitrary $\epsilon >0$, there exists a compact set $K$ such that for all $n$, $\mu^*(K^c),\mu_n(K^c) \leq \epsilon$. Since $g(y,y')$ is bounded, there exists a constant $C$ such that
    \begin{align}
          & \left|\iint g(y,y')d\mu_n(y) d\mu_n(y) - \iint g(y,y')d\mu^* d\mu^*\right| \nonumber            \\
        = & \left|\iint g(y,y') \mathbbm{1}_{K \times K}(d\mu_n d\mu_n - d\mu^* d\mu^*)\right| + C\epsilon.
    \end{align}
    Define
    \begin{align}
        G_n(y) = \int g(y,y') \mathbbm{1}_K d\mu_n(y'), \qquad G(y) = \int g(y,y') \mathbbm{1}_K d\mu^*(y').
    \end{align}
    We now require the kernel $g$ to have bounded derivative, and we have the following claims:
    \begin{claim} \label{claim_1}
        $G_n(y)$ is uniformly equicontinuous on $K$.
    \end{claim}
    \begin{proof}[Proof of Claim]
        By mean value theorem and condition \ref{output_kernel_conditions}, we get
        \begin{align}
                 & |G_n(y)-G_n(\tilde{y})| = \left|\int k(\norm{y-y'}) - k(\norm{\tilde{y}-y'}) \mathbbm{1}_K d\mu_n(y')\right| \nonumber                              \\
            \leq & \int \norm{y-\tilde{y}}\left|\frac{k(\norm{y-y'}) - k(\norm{\tilde{y}-y'})}{\norm{(y-y')-(\tilde{y}-y')}} \right|\mathbbm{1}_K d\mu_n(y') \nonumber \\
            =    & O(\norm{y-\tilde{y}}).
        \end{align}
    \end{proof}
    \begin{claim} \label{claim_2}
        $\sup_{y \in K}|G_n(y)-G(y)|=o(1)$.
    \end{claim}
    \begin{proof}[Proof of Claim]
        Since $K$ is compact, it suffices to show that for each $y \in K$, there exists $\delta>0$ such that
        \begin{align}
            \sup_{\tilde{y}\in B(y,\delta)} |G_n(\tilde{y})-G(\tilde{y})| \leq \epsilon.
        \end{align}
        Towards this end, by triangular inequality,
        \begin{align}
            |G_n(\tilde{y})-G(\tilde{y})| \leq |G_n(\tilde{y})-G_n(y)|+|G_n(y)-G(y)|+|G(y)-G(\tilde{y})|.
        \end{align}
        By the previous claim, the first term on the RHS vanishes because of equicontinuity. The second term vanishes because of weak convergence and the boundedness of the integrand $g(y,y')$. Finally, the last term vanishes because we can prove the equicontinuity of $G$ using a similar argument like the claim \ref{claim_1}.
    \end{proof}
    By claim \ref{claim_2} we have
    \begin{align}
        \iint g(y,y') \mathbbm{1}_{K \times K} d\mu_n d\mu_n = \iint g(y,y') \mathbbm{1}_{K \times K} d\mu^* d\mu_n +o(1) = \int G(y) \mathbbm{1}_K d\mu_n(y)+o(1).
    \end{align}
    Since $G$ is bounded and continuous by condition \ref{output_kernel_conditions}, using weak convergence of $\mu_n$ and we have
    \begin{align}
        \int G(y) \mathbbm{1}_K d\mu_n(y) = \int G(y) \mathbbm{1}_K d\mu^*n(y) + o(1) = \iint g(y,y')d\mu^* d\mu^* + o(1),
    \end{align}
    as desired.
\end{proof}
\begin{lemma}
    If $\mu^*$ has compact support in $Y$, then
    \begin{align}
        \iint p_n(x,x') \log(g(y,y')) d\mu_n d\mu_n = \iint p_n(x,x') \log(g(y,y')) d\mu^* d\mu^* +o(1).
    \end{align}
\end{lemma}
\begin{proof}
    By compactness, $\norm{y-y'}$ is bounded from above and $g(y,y')$ does not decay to zero. In that case, by lemma \ref{tsne:integrand_lemma_2}, we can replace the $p_n$ in the integral by $p$:
    \begin{align}
        \iint p_n(x,x') \log(g(y,y')) d\mu_n d\mu_n = \iint p(x,x') \log(g(y,y')) d\mu_n d\mu_n + o(1).
    \end{align}
    mimicking the proof of lemma \ref{entropy_lemma_1} we can show that
    \begin{align}
        \iint p(x,x') \log(g(y,y')) d\mu_n d\mu_n = \iint p(x,x') \log(g(y,y')) d\mu^* d\mu^* + o(1).
    \end{align}
    This finishes the proof.
\end{proof}
\begin{lemma}
    If $\mu^*$ has compact support in $Y$, then
    \begin{align}
        \sum_{i \neq j} p_{ij} \log \left(\frac{p_{ij}}{q_{ij}}\right) = \iint p(x,x') \log\left(\frac{p(x,x')}{q(y,y')}\right) d\mu^* d\mu^* +o(1)
    \end{align}
\end{lemma}
\begin{proof}
    This statement can be proved by combining the last 3 lemmas and lemma \ref{entropy_lemma_1}.
\end{proof}
Using the previous lemma and we can immediately get
\begin{proposition} \label{upper_bound_prop}
    If $\mu^*$ has compact support in $Y$, then
    \begin{align}
        \limsup_{n \to \infty} \inf_{Y \in (\mathbb{R}^s)^n} L_{n,\rho}(X,Y) \leq I_{\rho}(\mu^*).
    \end{align}
\end{proposition}
\subsection{Auxiliary Lemmas}\label{tsne:sec:extra_proofs}
This section contains lemmas that are necessary but are not a part of the flow
in the proofs of convergence. Recall that $\mu_n = \frac{1}{n}\sum_{i=1}^n
    \delta_{X_i,Y_i}$ is the empirical measure of both the input and output data
points. We just ignore the first or the second part based on the integration
variable because it is a probability measure. Let $Y = Y^* = (Y_1,\dots,Y_n)$
be a minimizer of $L$. Recall from (\ref{L_derivative}) that
\begin{align}
    \frac{\partial L_{n, \rho}(X,Y)}{\partial y_i} = -\sum_{j=1, j \neq i}^n (p_{ij}-q_{ij}) \frac{\frac{\partial}{\partial y_i}g(y_i,y_j)}{g(y_i,y_j)} = 0.
\end{align}
\begin{lemma} \label{appendix_lemma}
    Suppose that $\mu_X$ has either compact support and a $C^0$ density $f(x)$. Then for $\mu$-a.e. $x$ and $y$ we have
    \begin{align}
        \int (p(x,x')-q(y,y')) \frac{\frac{\partial}{\partial y}g(y,y')}{g(y,y')} d\mu(x',y') = 0.
    \end{align}
\end{lemma}
\begin{proof}
    Let $\alpha(x,y)$ be a test function. It suffices to show that
    \begin{align}
        \iint \alpha(x,y)(p(x,x')-q(y,y')) \frac{\frac{\partial}{\partial y}g(y,y')}{g(y,y')} d\mu(x',y') = 0.
    \end{align}
    Recalling the definition of $p_n$ and $q_n$, we see from (\ref{L_derivative}) that for all $i$,
    \begin{align}
        \frac{1}{n^2}\sum_{j=1,j \neq i}^n (p_n(X_i,X_j)-q_n(X_i,X_j)) \frac{\frac{\partial}{\partial Y_i}g(Y_i,Y_j)}{g(Y_i,Y_j)} = 0.
    \end{align}
    Multiply by $\alpha(X_i,Y_i)$ for each $i$ and the both sides are still zero. Afterwards, take the summation over $i$ and we get
    \begin{align}
          & \iint \alpha(x,y)(p_n(x,x')-q_n(y,y')) \frac{\frac{\partial}{\partial y}g(y,y')}{g(y,y')} d\mu_n d\mu_n \nonumber                \\
        = & \frac{1}{n^2}\sum_{1 \leq i,j \leq n} (p_n(X_i,X_j)-q_n(X_i,X_j)) \frac{\frac{\partial}{\partial Y_i}g(Y_i,Y_j)}{g(Y_i,Y_j)} = 0
    \end{align}
    if we require $k'(0) = 0$.
    By lemma \ref{tsne:integrand_lemma_2} and lemma \ref{output_kernel_int_lemma}, we have
    \begin{align}
        p_n(X_i,X_j)-q_n(Y_i,Y_j) = p(X_i,X_j)-q(Y_i,Y_j) + o(1).
    \end{align}
    Since both $\alpha$ and $\frac{\frac{\partial}{\partial y_i}g(y_i,y_j)}{g(y_i,y_j)}$ are bounded based on our assumptions, we have
    \begin{align}
          & \iint \alpha(x,y)(p_n(x,x')-q_n(y,y')) \frac{\frac{\partial}{\partial y}g(y,y')}{g(y,y')} d\mu_n d\mu_n \nonumber \\
        = & \iint \alpha(x,y)(p(x,x')-q(y,y')) \frac{\frac{\partial}{\partial y}g(y,y')}{g(y,y')} d\mu_n d\mu_n + o(1).
    \end{align}
    By lemma \ref{appendix_lemma_1} and lemma \ref{appendix_lemma_2} below we have
    \begin{align}
          & \iint \alpha(x,y)(p(x,x')-q(y,y')) \frac{\frac{\partial}{\partial y}g(y,y')}{g(y,y')} d\mu_n d\mu_n \nonumber \\
        = & \iint \alpha(x,y)(p(x,x')-q(y,y')) \frac{\frac{\partial}{\partial y}g(y,y')}{g(y,y')} d\mu d\mu + o(1).
    \end{align}
    Since $\alpha$ is an arbitrary test function,
    \begin{align}
        \int (p(x,x')-q(y,y')) \frac{\frac{\partial}{\partial y}g(y,y')}{g(y,y')} d\mu(x',y') = 0.
    \end{align}
\end{proof}
\begin{lemma} \label{appendix_lemma_1}
    The following holds as $n \to \infty$:
    \begin{align}
        \iint \alpha(x,y)p(x,x') \frac{\frac{\partial}{\partial y}g(y,y')}{g(y,y')} d\mu_n d\mu_n = \iint \alpha(x,y)p(x,x') \frac{\frac{\partial}{\partial y}g(y,y')}{g(y,y')} d\mu d\mu + o(1)
    \end{align}
\end{lemma}
\begin{lemma} \label{appendix_lemma_2}
    The following holds as $n \to \infty$:
    \begin{align}
        \iint \alpha(x,y)q(y,y') \frac{\frac{\partial}{\partial y}g(y,y')}{g(y,y')} d\mu_n d\mu_n = \iint \alpha(x,y)q(y,y') \frac{\frac{\partial}{\partial y}g(y,y')}{g(y,y')} d\mu d\mu + o(1)
    \end{align}
\end{lemma}
Both lemmas above can be proved via a similar strategy in the proof of lemma \ref{int_approx_lemma_2} by analyzing the diagonal terms.
\subsection{Proof of Theorem \ref{tsne:main_theorem_1}} \label{tsne:sec:main_theorem_1}
Write
\begin{align}
    d_{n,\rho}(X) = \inf_{Y \in (\mathbb{R}^s)^n}L_{n,\rho}(X,Y).
\end{align}
By proposition \ref{upper_bound_prop}, \ref{lower_bound_prop} and the fact that $\mu^* = \argmin_{\mu \in \tilde{\mathcal{P}}_X} I_{\rho}(\mu)$, there exists $\mu \in \tilde{\mathcal{P}}_X$ and a subsequence $n(k)$ such that
\begin{align}
    I_{\rho}(\mu^*) \geq \limsup_{n \to \infty} d_{n,\rho}(X) \geq \limsup_{n \to \infty} d_{n(k),\rho}(X) \geq I_{\rho}(\mu) \geq I_{\rho}(\mu^*).
\end{align}
Therefore the limit
\begin{align}
    \lim_{n \to \infty} d_{n(k),\rho}(X) = I_{\rho}(\mu) = I_{\rho}(\mu^*)
\end{align}
holds when we look at the sub-sequence $n(k)$ with the empirical measure $\mu$ converging to $\mu$ weakly. For the general limit, we can repeat the reasoning above by looking at any convergent sub-sequences of $d_{n,\rho}(X)$. We are able to find a sub-sequence inside each convergent sub-sequence and they all converges to $I_{\rho}(\mu^*)$. Therefore each convergent sub-sequence of $d_{n,\rho}(X)$ has the same limit and the theorem is proved. \qed
\subsection{Proof of Theorem \ref{tsne:main_theorem_2}}
We need another lemma before proving this theorem.
\begin{lemma} \label{home_stretch_lemma}
    For every $\mu \in \tilde{\mathcal{P}}_X$, there exists a function $h_{\mu}(x)$ that does not depend on $y$ such that
    \begin{align}
        \int p(x,x') \log(1+\norm{y-y'})d\mu = h_{\mu(x)} - \frac{\int g(y,y')d\mu}{\iint g(y,y')d\mu d\mu}.
    \end{align}
    Moreover, if $\mu_X$ has compact support, then $h_{\mu}$ is continuous on the support of $\mu_X$.
\end{lemma}
\begin{proof}
    Recall from lemma \ref{derivative_lemma} and lemma \ref{appendix_lemma}, we have
    \begin{align}
        0 = & \int (p(x,x')-q(y,y')) \frac{\frac{\partial}{\partial y}g(y,y')}{g(y,y')} d\mu(x',y') = \int p(x,x') \frac{\frac{\partial}{\partial y}g(y,y')}{g(y,y')} d \mu -  \frac{\int \frac{\partial}{\partial y}g(y,y')d\mu}{\iint g(y,y')d\mu d\mu} \nonumber \\
        =   & \nabla_y \left(\int p(x,x') \log (g(y,y')^{-1}) d \mu -  \frac{\int g(y,y')d\mu}{\iint g(y,y')d\mu d\mu}\right).
    \end{align}
    This means that there exists $h_{\mu}: \mathbb{R}^d \to \mathbb{R}$ such that
    \begin{align}
        h_{\mu(x)} = \int p(x,x') \log(g(y,y')^{-1})d\mu + \frac{\int g(y,y')d\mu}{\iint g(y,y')d\mu d\mu}.
    \end{align}
    Since $h_{\mu(x)}$ does not change with $y$, it suffices to show that $\int p(x,x') \log(g(y,y')^{-1})d\mu$ is continuous in $x$. Let $\epsilon>0$ and suppose $x_n \to x$. Since $p$ is continuous, and $\mu_X$ has compact support, we have for some constant $C$,
    \begin{align}
        \iint \log(g(y,y')^{-1})d\mu d\mu \leq C \int p(x,x') \log(g(y,y')^{-1})d\mu < \infty.
    \end{align}
    In particular, for almost every $y$, $\int \log(g(y,y')^{-1})d\mu < \infty$. Since $p(x,x')$ is continuous and $\mu_X$ has compact support, $p(x,x')$ is uniformly continuous. Specifically, there exists $\delta = \delta(y,\mu)$ such that if $\norm{x_n-x} \leq \delta$,
    \begin{align}
        \sup_x |p_n(x,x')-p(x,x')| \leq \frac{\epsilon}{\log(g(y,y')^{-1})d\mu}.
    \end{align}
    This means
    \begin{align}
        \left|\int(p_n(x,x')-p(x,x')) \log(g(y,y')^{-1})d \mu\right| < \epsilon,
    \end{align}
    which is the definition of continuity. The proof is finished.
\end{proof}
\begin{proof}[Proof of Theorem \ref{tsne:main_theorem_2}] \label{tsne:sec:main_theorem_2}
    We prove by contradiction. Define $D_n = \{\norm{y} \geq n\}$. If $\mu^*$ does not have compact support in $Y$ then $\mu(D_n)>0$ for all $n$. By lemma \ref{home_stretch_lemma}, there exists continuous $h(x)$ such that
    \begin{align}
        \int p(x,x') \log(g(y,y')^{-1})d\mu^* = h(x) -  \frac{\int g(y,y')d\mu^*}{\iint g(y,y')d\mu^* d\mu^*}.
    \end{align}
    Multiply both sides by $\mathbbm{1}_{D_n(x,y)}/\mu^*(D_n)$ and integrate with respect to the other variable:
    \begin{align}
         & \iint p(x,x') \log(g(y,y')^{-1}) \frac{\mathbbm{1}_{D_n(x,y)}}{\mu^*(D_n)}d\mu^* d\mu^* \nonumber                                                                                          \\
         & \qquad \qquad \qquad= \int h(x)\frac{\mathbbm{1}_{D_n(x,y)}}{\mu^*(D_n)}d\mu^* -  \frac{\iint g(y,y')\frac{\mathbbm{1}_{D_n(x,y)}}{\mu^*(D_n)} d\mu^* d\mu^*}{\iint g(y,y')d\mu^* d\mu^*}.
    \end{align}
    Since $h$ is continuous and $\mu_X$ has compact support, the first term of the RHS is bounded uniformly in $n$. Moreover, $g(y,y')$  is bounded so the second term is also bounded uniformly in $n$. It follows that the LHS is bounded uniformly in $n$. With this being said, because $p(x,x')>0$ and $\mu_X$ has compact support, there exists a constant $c$ such that $p(x,x') \geq c$ for all $x,x'$ in the support of $\mu_X$.

    Meanwhile, since the output kernel $g(y,y') = k(\norm{y-y'})$ is integrable, it
    has a decaying tail and there exists $C(n)$ such that $C(n) \to \infty$ and if
    $\norm{x-x'} \geq n/2$, $g(y,y')^{-1}=k(\norm{y-y'})^{-1} \geq log (1+C(n))$.
    We compute
    \begin{align}
             & \frac{1}{\mu^*(D_n)} \iint p(x,x') \log(g(y,y')^{-1}) \mathbbm{1}_{\norm{y} \geq n} d\mu^* d\mu^* \nonumber                                   \\
        \geq & \frac{1}{\mu^*(D_n)} \iint p(x,x') \log(g(y,y')^{-1}) \mathbbm{1}_{\norm{y} \geq n} \mathbbm{1}_{\norm{y-y'} \geq n/2}d\mu^* d\mu^* \nonumber \\
        \geq & c \frac{\log(1+C(n))}{\mu^*(D_n)} \iint \mathbbm{1}_{\norm{y} \geq n} \mathbbm{1}_{\norm{y'} \leq n/2}d\mu^* d\mu^* \nonumber                 \\
        =    & c\log(1+C(n)) \mu^*(\norm{Y'}\leq n/2).
    \end{align}
    Clearly,
    \begin{align}
        \lim_{n \to \infty} c\log(1+C(n)) \mu^*(\norm{Y'}\leq n/2) = \infty.
    \end{align}
    This contradicts with the boundedness. We must have that $\mu^*(D_n)=0$ for $n$ sufficiently large. That is, $\mu^*(\norm{Y} \geq n)=0$, so $\mu^*$ has compact support in $Y$.
\end{proof}
\bibliographystyle{hunsrtnat}
\setcitestyle{numbers}
\nocite{*}
\bibliography{citation}

\end{document}